\documentclass[11pt,a4paper]{coverletter}

\senderblock
  {Radheem Bin Razi}
  {Distributed Systems \& Cloud-Native Engineer}
  {Ilmenau, Germany \textbar\ \href{mailto:sheikh.radheem@gmail.com}{sheikh.radheem@gmail.com} \textbar\ \href{https://radheem.github.io/my-cv}{portfolio}}

\begin{document}

%% ===== ENGLISH =====
\recipient{Deutsche Bank}{Hiring Team}{}

\opening{Dear Hiring Team,}

\vspace{8pt}\noindent\textbf{1. Why Deutsche Bank?}\par\vspace{3pt}

Modernizing legacy financial infrastructure requires treating data pipelines and cloud transitions as continuous engineering challenges rather than one-off migrations. Deutsche Bank’s TDI strategy treats infrastructure modernization and AI integration as continuous engineering challenges rather than one-off migrations. That approach aligns with how I have spent the last five years designing distributed systems that scale quietly under real-world load. I am applying for the Technology, Data \& Innovation Internship 2026 to contribute to that transition while learning from Berlin’s engineering teams.

\vspace{8pt}\noindent\textbf{2. Why me?}\par\vspace{3pt}

- Built Go microservices with gRPC and NATS JetStream to replace a legacy Dapr stack, cutting message latency by 40\% and supporting hot-reloadable LLM tool integrations. - Consolidated a single-host 5G testbed into composable Docker Compose stacks, enabling remote gNBs to attach over shared bridges while streaming per-UE metrics to Kafka and InfluxDB. - Secured sensitive documents with client-side AES-256-GCM encryption and automated bidirectional Google Sheets synchronization to track application lifecycles without exposing credentials. I am available full-time immediately, open to part-time or working-student arrangements, and based in Germany with relocation readiness within two to three weeks; a qualifying offer would only require a standard contract since I am Blue Card-ready.

\closing{Sincerely,}{Radheem Bin Razi}

\clearpage
\selectlanguage{ngerman}

%% ===== DEUTSCH =====
\recipient{Deutsche Bank}{Personalabteilung}{}

\opening{Sehr geehrtes Hiring-Team,}

\vspace{8pt}\noindent\textbf{1. Warum die Deutsche Bank?}\par\vspace{3pt}

Die Modernisierung veralteter Finanzinfrastrukturen erfordert es, Datenpipelines und Cloud-Migrationen als kontinuierliche Engineering-Herausforderungen und nicht als einmalige Projekte zu betrachten. Die TDI-Strategie der Deutschen Bank betrachtet die Modernisierung der Infrastruktur und die KI-Integration ebenfalls als kontinuierliche Engineering-Herausforderungen statt als einmalige Migrationen. Dieser Ansatz entspricht genau der Art und Weise, wie ich in den letzten fünf Jahren verteilte Systeme entworfen habe, die unter realer Last unauffällig skalieren. Ich bewerbe mich für das Technology, Data \& Innovation Internship 2026, um an dieser Transformation mitzuwirken und gleichzeitig von den Engineering-Teams in Berlin zu lernen.

\vspace{8pt}\noindent\textbf{2. Warum ich?}\par\vspace{3pt}

- Entwicklung von Go-Mikroservices mit gRPC und NATS JetStream zum Ersatz eines veralteten Dapr-Stacks, wodurch die Nachrichtenlatenz um 40 \% reduziert und hot-reloadable LLM-Tool-Integrationen unterstützt wurden. - Konsolidierung eines Single-Host-5G-Testbeds zu kompositionsfähigen Docker Compose Stacks, wodurch entfernte gNBs über Shared Bridges angebunden werden konnten und Per-UE-Metriken parallel an Kafka und InfluxDB gestreamt wurden. - Verschlüsselung sensibler Dokumente auf Client-Seite mittels AES-256-GCM sowie automatisierte bidirektionale Google Sheets-Synchronisation zur Nachverfolgung von Application Lifecycles, ohne dass Credentials offengelegt werden. Ich stehe ab sofort vollzeitlich zur Verfügung und bin auch für Teilzeit- oder Werkstudentenmodelle offen. Ich befinde mich bereits in Deutschland und bin innerhalb von zwei bis drei Wochen umzugsbereit; ein entsprechendes Angebot würde lediglich einen Standardarbeitsvertrag erfordern, da ich Blue Card-ready bin.

\closing{Mit freundlichen Grüßen,}{Radheem Bin Razi}

\end{document}
